\documentclass[../../../../SoftwareRequirementsSpecification.tex]{subfiles}
\begin{document}
% Richiede le macro definite in src/macros/useCaseMacros.tex
% (incluse nel preambolo di SoftwareRequirementsSpecification.tex)

\ucstart
\begin{xltabular}{\textwidth}{|l|X|}
  \hline
  \ucid{PT-09} \\ \hline
  \ucname{Personalizzazione di una Scheda Assegnata} \\ \hline
  \ucactor{PT} \\ \hline
  \ucdescription{Il PT modifica le sessioni di allenamento di una
    scheda che un atleta sta già seguendo. La modifica vale solo per
    quell'atleta.} \\ \hline
  \ucprecond{Il PT deve essere autenticato. L'atleta scelto deve
    avere almeno una scheda assegnata in corso, con assegnazione
    confermata: le bozze non sono personalizzabili (vedi Nota 5 del
    caso d'uso PT-05).} \\ \hline
  \ucpostcond{Le sessioni di allenamento prescritte a quell'atleta per
    la settimana scelta sono modificate. La scheda non cambia. Le
    sessioni di allenamento degli altri atleti che seguono la scheda
    non cambiano.} \\ \hline
  \ucmainflow{
    \item Il PT si trova sulla pagina di dettaglio di un atleta,
      raggiunta tramite il caso d'uso PT-08.
    \item Il PT preme su una scheda dell'elenco delle schede
      assegnate.
    \item Il sistema mostra il dettaglio dell'assegnazione: nome
      della scheda, periodo di validità ed elenco delle settimane di
      allenamento. Il sistema evidenzia la settimana corrente e
      segnala le settimane che contengono personalizzazioni.
    \item Il PT preme su una settimana, scelta tra la settimana
      corrente e quelle successive (vedi Nota 2).
    \item Il sistema mostra le sessioni di allenamento della
      settimana con i valori previsti per quell'atleta (vedi
      Nota 1).
    \item Il PT preme su un esercizio.
    \item Il sistema mostra il form con i valori correnti
      dell'esercizio:
      \begin{itemize}
        \item peso;
        \item numero di serie/ripetizioni;
        \item tempo di recupero tra una serie e l'altra (in
          secondi).
      \end{itemize}
    \item Il PT modifica i valori che vuole cambiare.
    \item Il PT preme il pulsante "Salva".
    \item Il sistema registra la personalizzazione per la settimana
      scelta.
    \item Il sistema invia un'email di notifica all'Atleta per
      informarlo che la sua scheda è cambiata.
  } \\ \hline
  \ucaltflow{
    \textbf{2a. Nessuna scheda assegnata} \newline
    - L'atleta non ha schede assegnate. Il sistema mostra la
    sezione vuota. \newline
    - Il flusso termina. \newline
    \textbf{4a. Scheda terminata} \newline
    - Tutte le settimane della scheda sono trascorse: nessuna
    settimana è modificabile. \newline
    - Il sistema mostra un messaggio ("Scheda terminata: non ci sono
    settimane modificabili"). \newline
    - Il flusso termina. \newline
    \textbf{7a. Ripristino del valore originale} \newline
    - Se l'esercizio è personalizzato, il sistema mostra
    accanto ai valori il pulsante "Ripristina". \newline
    - Il PT preme il pulsante. Il sistema riporta l'esercizio ai valori
    previsti dalla scheda. \newline
    - Il flusso riprende dal passo 5.
  } \\ \hline
  \ucnote{
    \textbf{Nota 1:} quali esercizi l'atleta svolge è deciso dalle
    sessioni di allenamento configurate nel blocco mensile (PT-04) e
    non è modificabile qui. Sono invece memorizzati sull'assegnazione,
    per ciascun esercizio di ciascuna sessione: il carico, che il PT
    inserisce al momento dell'assegnazione (PT-05), e le
    serie/ripetizioni e il tempo di recupero, che restano quelli
    della scheda finché il PT non li modifica con questo caso d'uso.
    La modifica agisce sull'assegnazione e non sulla scheda: la
    scheda resta invariata e gli altri atleti che la seguono non sono
    toccati.\newline
    \textbf{Nota 2:} sono modificabili solo la settimana corrente e
    le settimane successive. Le settimane già trascorse sono in sola
    lettura, perché l'atleta si è già allenato su quanto vi era
    prescritto.\newline
    \textbf{Nota 3:} lo statement richiede che la tipologia di
    esercizi in una scheda vari solo da mese a mese. Questo caso
    d'uso rispetta il vincolo: la personalizzazione tocca solo peso,
    serie/ripetizioni e tempo di recupero, che lo statement lascia
    liberi di variare da una sessione di allenamento all'altra. Quali
    esercizi l'atleta svolge resta fissato per l'intero blocco mensile
    dalle sessioni di allenamento configurate in PT-04.\newline
    \textbf{Nota 4:} la personalizzazione non permette di: aggiungere,
    sostituire o togliere esercizi, cambiare il numero o la
    composizione delle sessioni, cambiare la durata della scheda.
    Questi elementi sono parte della scheda e valgono per tutti gli
    atleti che la seguono. Per interventi di questo tipo, il PT deve
    duplicare la scheda e assegnare la copia.
  } \\ \hline
\end{xltabular}
\end{document}
